\definecolor{links}{HTML}{2A1B81}
\hypersetup{colorlinks,linkcolor=,urlcolor=links}

\usetheme{Boadilla}
\usecolortheme{seahorse}
\usefonttheme{serif}
\beamertemplatenavigationsymbolsempty

\usepackage{luacode}
\usepackage{luatexja}
\usepackage{pgfpages}
\usepackage[osf]{mathpazo}

\begin{luacode*}
  USE_IPAFONT = os.getenv"USE_IPAFONT"
  USE_YUFONT = os.getenv"USE_YUFONT"
  
  if USE_YUFONT == "true" then
    tex.sprint("\\AtBeginDocument{\\usepackage[yu-osx, deluxe, expert]{luatexja-preset}}")
    tex.sprint("\\AtBeginDocument{\\setmainjfont[BoldFont=YuGothic Bold]{YuGothic Medium}}")
  elseif USE_IPAFONT == "true" then
    tex.sprint("\\AtBeginDocument{\\usepackage[ipaex, deluxe, expert]{luatexja-preset}}")
  else
    tex.sprint("\\AtBeginDocument{\\usepackage[hiragino-pro, deluxe, expert]{luatexja-preset}}")
    tex.sprint("\\AtBeginDocument{\\setmainjfont[BoldFont=Hiragino Kaku Gothic Pro W6]{Hiragino Kaku Gothic Pro W3}}")
  end
\end{luacode*}

\usepackage{epigraph}
\usepackage{etoolbox}
\usepackage{tikz}
\usepackage{framed}
\usepackage{libertine}
\usepackage[final]{listings}
\usepackage{amsmath}
\usepackage{mathtools}

%\setbeameroption{show notes on second screen=right}

%\usetikzlibrary{arrows,automata,shapes,backgrounds}

\setmainfont[Numbers=OldStyle, BoldFont=Palatino Bold]{Palatino}
\setsansfont{CMU Sans Serif}
\setmonofont{CMU Typewriter Text}

\title[Mental Jinro]{%
  \href{https://github.com/y-yu/mental-jinro-slide}{Mental Jinro} \\
  \href{http://connpass.com/event/30661/}{\normalsize tsukuba.pm \#3}
}
\author{吉村 優}
\date{May 14, 2016}
\institute[\url{https://twitter.com/\_yyu\_}]{%
  \url{https://twitter.com/\_yyu\_}\\
  \url{http://qiita.com/yyu}\\
  \url{https://github.com/y-yu}\\
}

\newfontfamily\quotefont[Ligatures=TeX]{Linux Libertine O} % selects Libertine as the quote font

\newcommand*\quotesize{60} % if quote size changes, need a way to make shifts relative
% Make commands for the quotes
\newcommand*{\openquote}
   {\tikz[remember picture,overlay,xshift=-4ex,yshift=-2.5ex]
   \node (OQ) {\quotefont\fontsize{\quotesize}{\quotesize}\selectfont``};\kern0pt}

\newcommand*{\closequote}[1]
  {\tikz[remember picture,overlay,xshift=1.5ex,yshift={#1}]
   \node (CQ) {\quotefont\fontsize{\quotesize}{\quotesize}\selectfont''};}

\newcommand*\shadedauthorformat{} % define format for the author argument

% Now a command to allow left, right and centre alignment of the author
\newcommand*\authoralign[1]{%
  \if#1l
    \def\authorfill{}\def\quotefill{\hfill}
  \else
    \if#1r
      \def\authorfill{\hfill}\def\quotefill{}
    \else
      \if#1c
        \gdef\authorfill{\hfill}\def\quotefill{\hfill}
      \else\typeout{Invalid option}
      \fi
    \fi
  \fi}
% wrap everything in its own environment which takes one argument (author) and one optional argument
% specifying the alignment [l, r or c]
%
\newenvironment{shadequote}[2][l]%
{\authoralign{#1}
\ifblank{#2}
   {\def\shadequoteauthor{}\def\yshift{-2ex}\def\quotefill{\hfill}}
   {\def\shadequoteauthor{\par\authorfill\shadedauthorformat{#2}}\def\yshift{2ex}}
\begin{quote}\openquote}
{\shadequoteauthor\quotefill\closequote{\yshift}\end{quote}}

\newfontfamily\listingfont{Menlo}

\lstset{
  basicstyle=\listingfont
}

\lstdefinestyle{cpp}{
  basicstyle=\listingfont\footnotesize,        % the size of the fonts that are used for the code
  breakatwhitespace=false,         % sets if automatic breaks should only happen at whitespace
  language=C++,
  firstnumber=0,
  xleftmargin=2em,
  framexleftmargin=1.5em,
  numberstyle=\listingfont\scriptsize,
  breaklines=true,                 % sets automatic line breaking
  captionpos=b,                    % sets the caption-position to bottom
  commentstyle=\listingfont\color{green},    % comment style
  extendedchars=true,              % lets you use non-ASCII characters; for 8-bits encodings only, does not work with UTF-8
  keepspaces=true,                 % keeps spaces in text, useful for keeping indentation of code (possibly needs columns=flexible)
  keywordstyle=\listingfont\color{blue},       % keyword style
  rulecolor=\listingfont\color{black},         % if not set, the frame-color may be changed on line-breaks within not-black text (e.g. comments (green here))
  showspaces=false,                % show spaces everywhere adding particular underscores; it overrides 'showstringspaces'
  showstringspaces=false,          % underline spaces within strings only
  showtabs=false,                  % show tabs within strings adding particular underscores
  showlines=false,
  stringstyle=\listingfont\color{gray},        % string literal style
  tabsize=2,	                   % sets default tabsize to 2 spaces
  title=\lstname                   % show the filename of files included with \lstinputlisting; also try caption instead of title
}

\lstdefinestyle{vm}{
  basicstyle=\listingfont\footnotesize,        % the size of the fonts that are used for the code
  breakatwhitespace=false,         % sets if automatic breaks should only happen at whitespace
  morekeywords={char,split,jmp,match},
  numbers=left,
  firstnumber=0,
  xleftmargin=2em,
  framexleftmargin=1.5em,
  numberstyle=\listingfont\scriptsize,
  breaklines=true,                 % sets automatic line breaking
  captionpos=b,                    % sets the caption-position to bottom
  commentstyle=\listingfont\color{green},    % comment style
  extendedchars=true,              % lets you use non-ASCII characters; for 8-bits encodings only, does not work with UTF-8
  keepspaces=true,                 % keeps spaces in text, useful for keeping indentation of code (possibly needs columns=flexible)
  keywordstyle=\listingfont\color{blue},       % keyword style
  rulecolor=\listingfont\color{black},         % if not set, the frame-color may be changed on line-breaks within not-black text (e.g. comments (green here))
  showspaces=false,                % show spaces everywhere adding particular underscores; it overrides 'showstringspaces'
  showstringspaces=false,          % underline spaces within strings only
  showtabs=false,                  % show tabs within strings adding particular underscores
  showlines=false,
  stringstyle=\listingfont\color{gray},        % string literal style
  tabsize=2,	                   % sets default tabsize to 2 spaces
  title=\lstname                   % show the filename of files included with \lstinputlisting; also try caption instead of title
}

\lstdefinestyle{grass}{
  basicstyle=\listingfont\scriptsize,        % the size of the fonts that are used for the code
  breakatwhitespace=false,         % sets if automatic breaks should only happen at whitespace
  xleftmargin=2em,
  framexleftmargin=1.5em,
  numberstyle=\listingfont\scriptsize,
  breaklines=true,                 % sets automatic line breaking
  captionpos=b,                    % sets the caption-position to bottom
  commentstyle=\listingfont\color{green},    % comment style
  extendedchars=true,              % lets you use non-ASCII characters; for 8-bits encodings only, does not work with UTF-8
  keepspaces=true,                 % keeps spaces in text, useful for keeping indentation of code (possibly needs columns=flexible)
  keywordstyle=\listingfont\color{blue},       % keyword style
  rulecolor=\listingfont\color{black},         % if not set, the frame-color may be changed on line-breaks within not-black text (e.g. comments (green here))
  showspaces=false,                % show spaces everywhere adding particular underscores; it overrides 'showstringspaces'
  showstringspaces=false,          % underline spaces within strings only
  showtabs=false,                  % show tabs within strings adding particular underscores
  showlines=false,
  stringstyle=\listingfont\color{gray},        % string literal style
  tabsize=2,	                   % sets default tabsize to 2 spaces
  title=\lstname                   % show the filename of files included with \lstinputlisting; also try caption instead of title
}


\makeatletter
\def\@fnsymbol#1{\ensuremath{\ifcase#1\or *\or \dagger\or \ddagger\or
   \mathsection\or \mathparagraph\or \|\or **\or \dagger\dagger
   \or \ddagger\ddagger \else\@ctrerr\fi}}
\makeatother

\renewcommand{\thefootnote}{\fnsymbol{footnote}}

\makeatletter
\def\Ldescription{%
  \@ifnextchar[{\beamer@testforospec}{\beamer@descdefault\beamer@descriptionwidth\@@Ldescription}%
}

\def\beamer@testforospec[{\@ifnextchar<{\beamer@scandefaultospec[}{\@Ldescription[}}%

\def\beamer@scandefaultospec[#1]{\def\beamer@defaultospec{#1}\Ldescription}

\def\@Ldescription[#1]{%
\setbox\beamer@tempbox=\hbox{\def\insertdescriptionitem{#1}
  \usebeamertemplate**{description item}}%
\beamer@descdefault\wd\beamer@tempbox\@@description%
}%

\def\@@Ldescription{%
  \beamer@descdefault35pt%
  \list
  {}
  {\labelwidth\beamer@descdefault\leftmargin2.8em\let\makelabel\beamer@Ldescriptionitem}%
  \beamer@cramped%
  \raggedright
  \beamer@firstlineitemizeunskip%
}

\def\endLdescription{\ifhmode\unskip\fi\endlist}
\long\def\beamer@Ldescriptionitem#1{%
  \def\insertdescriptionitem{#1}%
  \hspace\labelsep{\parbox[b]{\dimexpr\textwidth-\labelsep\relax}{%
        \usebeamertemplate**{description item}%
    }}}
\makeatother

\usetikzlibrary{shapes.callouts} 

\pgfkeys{%
    /calloutquote/.cd,
    width/.code                   =  {\def\calloutquotewidth{#1}},
    position/.code                =  {\def\calloutquotepos{#1}}, 
    author/.code                  =  {\def\calloutquoteauthor{#1}},
    /calloutquote/.unknown/.code   =  {\let\searchname=\pgfkeyscurrentname
                                 \pgfkeysalso{\searchname/.try=#1,                                
    /tikz/\searchname/.retry=#1},\pgfkeysalso{\searchname/.try=#1,
                                  /pgf/\searchname/.retry=#1}}
                            }  


\newcommand\calloutquote[2][]{%
       \pgfkeys{/calloutquote/.cd,
         width               = 5cm,
         position            = {(0,-1)},
         author              = {}}
  \pgfqkeys{/calloutquote}{#1}                   
  \node [rectangle callout,callout relative pointer={\calloutquotepos},text width=\calloutquotewidth,/calloutquote/.cd,
     #1] (tmpcall) at (0,0) {\hfil#2\hfil};
  \node at (tmpcall.pointer){\calloutquoteauthor};    
}


\newcommand\ballref[1]{%
\tikz \node[circle, shade,ball color=structure.fg,inner sep=0pt,%
  text width=8pt,font=\tiny,align=center] {\color{white}\ref{#1}};
}

%\everymath{\displaystyle}

\begin{document}

\frame{\maketitle}

\section{自己紹介}
\begin{frame}[fragile]
  \frametitle{自己紹介}
  
  \begin{columns}
    \begin{column}{0.4\textwidth}
      \centering
      \begin{figure}
        \includegraphics[width=0.9\textwidth]{img/bird2x.png}
      \end{figure}
    \end{column}
    \begin{column}{0.6\textwidth}
      \begin{itemize}
        \item<2-> \href{https://www.coins.tsukuba.ac.jp/}{筑波大学 情報科学類} 学士（COINS11）
        \item<3-> \href{http://www.word-ac.net/}{WORD編集部} OB
        \item<4-> \href{http://logic.cs.tsukuba.ac.jp/index-j.html}{プログラム論理研究室} OB
      \end{itemize}
    \end{column}
  \end{columns}
\end{frame}

\section{``Mental Jinro''とは？}
\begin{frame}[fragile]
  \frametitle{``Mental Jinro''とは？}
  
   \uncover<2->{
    \begin{block}{}
      \begin{shadequote}[r]{}
        \textbf{Mental Jinro}は\textbf{人狼}から\textbf{ゲームマスター}を排除したゲームである。
      \end{shadequote}
    \end{block}
  }

  \begin{center}
    \uncover<3->{
      \begin{tikzpicture}
        \calloutquote[author={},width=6cm,position={(-1.2,-0.5)},fill=green!30,rounded corners]{\textbf{ゲームマスター}とは何か？}
      \end{tikzpicture}
    }

    \uncover<4->{
      \begin{tikzpicture}
        \calloutquote[author={},width=6cm,position={(1,-0.5)},fill=blue!30,rounded corners]{そもそも\textbf{人狼}とは何か？}
      \end{tikzpicture}
    }
  \end{center}
\end{frame}

\subsection{人狼とは？}
\begin{frame}[fragile]
  \frametitle{人狼とは？}

  \begin{exampleblock}{人狼\footnote[frame]{人狼には様々なルールがあるが、このスライドではこのルールを用いる}}
    \begin{shadequote}[r]{\scriptsize \href{https://ja.wikipedia.org/wiki/\%E6\%B1\%9D\%E3\%81\%AF\%E4\%BA\%BA\%E7\%8B\%BC\%E3\%81\%AA\%E3\%82\%8A\%E3\%82\%84\%3F}{Wikipedia ------ 汝は人狼なりや？}}
      \begin{itemize}
        \item<2-> プレイヤーはそれぞれが\textbf{村人}と村人に化けた\textbf{人狼}となり、自分自身の正体がばれないように他のプレイヤーと交渉して正体を探る
        \item<3-> ゲームは半日単位で進行し、昼には全プレイヤーの投票により決まった人狼容疑者の\textbf{処刑}が、夜には人狼による村人の\textbf{襲撃}が行われる
        \item
          \begin{itemize}
            \item<4-> 全ての人狼を処刑することができれば村人チームの勝ち
            \item<5-> 人狼と同じ数まで村人を減らすことができれば人狼チームの勝ち
          \end{itemize}
      \end{itemize}
    \end{shadequote}
  \end{exampleblock}
\end{frame}

\subsection{人狼の役職}
\begin{frame}[fragile]
  \frametitle{人狼の役職}

  \begin{block}{人狼に必要な役職}
    参加者を次の役職に分ける必要がある
    \begin{itemize}
      \item<2-> 村人
      \item<3-> 人狼
      \item<4-> \textbf{ゲームマスター}
    \end{itemize}
  \end{block}

  \begin{center}
    \uncover<5->{
      \begin{tikzpicture}
        \calloutquote[author={},width=6cm,position={(-1,-0.5)},fill=green!30,rounded corners]{\textbf{ゲームマスター}とは何か？}
      \end{tikzpicture}
    }
  \end{center}
\end{frame}

\subsection{ゲームマスターと公平性}
\begin{frame}[fragile]
  \frametitle{ゲームマスターと公平性}

  \begin{block}{ゲームマスターの役割}
    \begin{itemize}
      \item<2-> 人狼と村人のチーム分けをする
      \item<3-> 人狼に襲撃された村人を村人チームに宣告する
      \item<4-> 人狼と村人の数を管理し、どちらかのチームが勝利した時それを宣言する
    \end{itemize}
  \end{block}

  \begin{center}
    \uncover<5->{
      \begin{tikzpicture}
        \calloutquote[author={},width=5cm,position={(-1.3,-0.5)},fill=green!30,rounded corners]{\textbf{ゲームマスター}とは審判}
      \end{tikzpicture}
    }

    \uncover<6->{
      \begin{tikzpicture}
        \calloutquote[author={},width=7cm,position={(1.2,-0.5)},fill=blue!30,rounded corners]{ゲームマスターが\textbf{不公平}だったら？}
      \end{tikzpicture}
    }

    \uncover<7->{
      \begin{tikzpicture}
        \calloutquote[author={},width=2cm,position={(-1,-0.5)},fill=red!30,rounded corners]{\textbf{大問題！}}
      \end{tikzpicture}
    }
  \end{center}
\end{frame}

\subsection{Mental Jinro}
\begin{frame}[fragile]
  \frametitle{Mental Jinro}

  \begin{center}
    \begin{tikzpicture}
      \calloutquote[author={\includegraphics[width=1cm]{./img/bird2x.png}},width=10cm,position={(-1.1,-0.5)},fill=green!30,rounded corners]{\huge\textbf{ゲームマスター}を消そう！}
    \end{tikzpicture}
  \end{center}

  \uncover<2->{
    \begin{alertblock}{ゲームマスターが消えると……}
      \begin{itemize}
        \item<3-> チーム分けはどうする？
        \item<4-> 襲撃された村人の情報をどう伝える？
        \item<5-> 勝敗は誰が判断する？ 
      \end{itemize}
    \end{alertblock}
  }
\end{frame}

\begin{frame}
  \centering
  {\Huge Mental Jinroを支える技術 \\ \uncover<2->{\textbf{コミットメント}}}
\end{frame}

\section{コミットメント}

\subsection{コイントスゲーム}
\begin{frame}[fragile]
  \frametitle{コイントスゲーム}

  \begin{exampleblock}{}
    アリスとボブの二人がいるとする
    \begin{enumerate}
      \item<2-> アリスがコインの``表''または``裏''を紙に書き、紙を封筒に入れる
      \item<3-> ボブはコインを投げる
      \item<4-> 封筒から紙を取り出し、
        \begin{itemize}
          \item アリスの予想とコインの結果が同じなら、アリスの勝利
          \item アリスの予想とコインの結果が違えば、ボブの勝利
        \end{itemize}
      \item<5-> このゲームは\textbf{電話上}で行う
    \end{enumerate}
  \end{exampleblock}

  \begin{center}
    \uncover<6->{
      \begin{tikzpicture}
        \calloutquote[author={},width=6cm,position={(-1.2,-0.3)},fill=green!30,rounded corners]{アリスが予想を\textbf{反故}にする？}
      \end{tikzpicture}
    }

    \uncover<7->{
      \begin{tikzpicture}
        \calloutquote[author={},width=6.5cm,position={(1.1,-0.3)},fill=blue!30,rounded corners]{ボブがアリスの主張を信じない？}
      \end{tikzpicture}
    }
  \end{center}
\end{frame}

\begin{frame}[fragile]
  \frametitle{コイントスゲーム}

  \begin{block}{プロトコル}
    \begin{enumerate}
      \item<2-> ボブは$p = 2q + 1$となる大きな素数$p, q$をランダムに生成して、
        $\mathbb{Z}^*_p$\footnote[frame]{整数$x \bmod p$かつ$xy \equiv 1 \pmod{p}$となる逆元$y$が存在する$x$の集合である}
        の位数$q$の部分群$G$から生成元$g, v$をランダムに選択して$p, q, g, v$をアリスへ送信する
      \item<3-> アリスは$p, q, g, v$を検証し、表と予想するなら$m := 1$を選択し、裏と予想するなら$m := q - 1$を選択し、
        乱数$r \in \{1, \dots, q - 1\}$を用いて$c := g^r v^m \bmod p$計算し$c$をボブへ送信する
      \item<4-> ボブはコイントスをして、結果をアリスへ送信する
      \item<5-> アリスは$r, m$を公開する
      \item<6-> ボブは$c \equiv g^r v^m \pmod{n}$を検証する
    \end{enumerate}
  \end{block}
\end{frame}

\begin{frame}{目次}
  \tableofcontents
\end{frame}

\end{document}
